\begin{table}[H]
\centering
\caption{Nama Daerah Perbandingan Cluster STATIS Tahunan vs IPM (UNDP) (2023)} \label{tab:app2statis_ipm_name_2023}
\begin{adjustbox}{max width=\textwidth}
\begin{tabular}{|>{\centering\arraybackslash}p{4cm}|p{10cm}|}
\hline
\multicolumn{1}{|c|}{\textbf{Kategori}} & \multicolumn{1}{c|}{\textbf{Daftar Wilayah}} \\
\hline
Cl. 1 (Rendah) \& Sangat Tinggi & Kota Jakarta Selatan, Kota Jakarta Timur, Kota Jakarta Pusat, Kota Jakarta Barat, Kota Jakarta Utara, Kota Bandung, Kota Bekasi, Kota Depok, Kota Tangerang Selatan \\
\hline
Cl. 1 (Rendah) \& Tinggi & Bekasi, Kota Bogor, Kota Sukabumi, Kota Cimahi, Kota Tangerang \\
\hline
Cl. 2 (Sedang) \& Tinggi & Kep. Seribu, Bogor, Bandung, Ciamis, Kuningan, Cirebon, Sumedang, Subang, Purwakarta, Karawang, Kota Cirebon, Kota Tasikmalaya, Kota Banjar, Tangerang, Kota Cilegon, Kota Serang \\
\hline
Cl. 2 (Sedang) \& Sedang & Sukabumi, Garut, Indramayu, Bandung Barat \\
\hline
Cl. 3 (Tinggi) \& Sedang & Cianjur, Tasikmalaya, Majalengka, Pangandaran, Pandeglang, Lebak, Serang \\
\hline
\end{tabular}
\end{adjustbox}
\end{table}